% Qiushi Article Template
% Compile with LuaLaTeX.
% If the LuaLaTeX font cache is unavailable, prefix the compile command with TEXMFVAR=/tmp/qiushi_texmf_var.
% Keep figure files under figures/ or ../../workspace/figures/.
% Use refs.bib with \cite{...} when references are used.
% LaTeX supplies section numbers; section titles contain title text only.
% ctexart two-column mode: figure is one-column and figure* is full-width.
\documentclass[UTF8,a4paper,10pt,twocolumn]{ctexart}

\usepackage[
  a4paper,
  top=16mm,
  bottom=18mm,
  left=15mm,
  right=15mm,
  columnsep=7.2mm
]{geometry}
\usepackage{graphicx}
\usepackage{amsmath,amssymb,mathtools,bm}
\usepackage{amsthm}
\usepackage{siunitx}
\usepackage{booktabs}
\usepackage{tabularx}
\usepackage{array}
\usepackage{xurl}
\usepackage{xcolor}
\usepackage{abstract}
\usepackage[numbers,sort&compress]{natbib}
\usepackage[colorlinks=true,linkcolor=blue!55!black,citecolor=blue!55!black,urlcolor=blue!55!black]{hyperref}

\graphicspath{{figures/}{../figures/}{./}{../../workspace/figures/}}
\sisetup{detect-all=true,per-mode=symbol}
\Urlmuskip=0mu plus 2mu
\emergencystretch=1.8em
\raggedbottom

\setlength{\parindent}{2em}
\setlength{\parskip}{0pt}
\setlength{\textfloatsep}{8pt plus 2pt minus 2pt}
\setlength{\floatsep}{7pt plus 2pt minus 2pt}
\setlength{\intextsep}{7pt plus 2pt minus 2pt}
\setlength{\bibsep}{2pt plus 0.2ex}

\ctexset{
  section = {
    format = {\raggedright\bfseries\large},
    beforeskip = 10pt plus 2pt minus 1pt,
    afterskip = 4pt
  },
  subsection = {
    format = {\raggedright\bfseries\normalsize},
    beforeskip = 8pt plus 1pt minus 1pt,
    afterskip = 3pt
  }
}

\newcolumntype{Y}{>{\raggedright\arraybackslash\hspace{0pt}}X}
\newcolumntype{L}[1]{>{\raggedright\arraybackslash\hspace{0pt}}p{#1}}

\providecommand{\dd}{\mathop{}\!\mathrm{d}}
\providecommand{\dv}[2]{\frac{\dd #1}{\dd #2}}
\providecommand{\pdv}[2]{\frac{\partial #1}{\partial #2}}
\providecommand{\vect}[1]{\boldsymbol{#1}}
\providecommand{\mat}[1]{\mathbf{#1}}
\providecommand{\E}{\mathbb{E}}
\providecommand{\Var}{\operatorname{Var}}
\providecommand{\argmin}{\operatorname*{arg\,min}}

\newtheorem{theorem}{定理}
\newtheorem{proposition}{命题}
\newtheorem{lemma}{引理}
\newtheorem{corollary}{推论}
\theoremstyle{definition}
\newtheorem{definition}{定义}
\theoremstyle{remark}
\newtheorem{remark}{注}

\title{概括主结果的中文题名}
\author{作者姓名\\\small 作者单位}
\date{\today}

\begin{document}

\twocolumn[
\maketitle

\begin{onecolabstract}
摘要正文。
\end{onecolabstract}
\vspace{0.7em}
]

\section{引言}

正文。

\section{方法}

正文。

\begin{equation}
  \mathcal{L}u = f .
  \label{eq:example-zh}
\end{equation}

\begin{figure}[t]
  \centering
  \fbox{\rule{0pt}{1.35in}\rule{0.90\columnwidth}{0pt}}
  \caption{图注。\label{fig:example-zh}}
\end{figure}

\section{结果}

正文。

\begin{figure*}[t]
  \centering
  \fbox{\rule{0pt}{1.55in}\rule{0.96\textwidth}{0pt}}
  \caption{需要跨双栏保持图内文字、面板和比较关系清晰的高信息密度图。\label{fig:two-column-zh}}
\end{figure*}

\begin{equation}
  \int_{\Omega} A(x)\,B(x)\,\dd x
  =
  \int_{\Omega} C(x)\,\dd x .
  \label{eq:wide-example-zh}
\end{equation}

\begin{table}[b]
  \caption{表注。\label{tab:example-zh}}
  \begin{tabular}{lcc}
  \toprule
  行 & 列 A & 列 B \\
  \midrule
  A & -- & -- \\
  B & -- & -- \\
  \bottomrule
  \end{tabular}
\end{table}

\begin{table*}[t]
  \caption{自然宽度和信息密度需要跨双栏呈现的比较表。\label{tab:two-column-zh}}
  \begin{tabular*}{\textwidth}{@{\extracolsep{\fill}}lccccc@{}}
  \toprule
  方法 & 假设 & 观测量 & 对照 & 不确定性 & 结果 \\
  \midrule
  A & -- & -- & -- & -- & -- \\
  B & -- & -- & -- & -- & -- \\
  \bottomrule
  \end{tabular*}
\end{table*}

% 双栏稿件仅在对象按最终字号无法保持单栏可读性时使用 figure* 或
% table*。公式优先在语义断点使用 aligned、split 或 multline；不得用
% resizebox 或不可读的小字号掩盖过宽对象。
% 必要时的跨栏语法：
% \begin{figure*}[t]\centering\includegraphics[width=\textwidth]{...}
%   \caption{...}\end{figure*}
% \begin{table*}[t]\centering\caption{...}\begin{tabular}{...}
%   ...\end{tabular}\end{table*}
% 语义断行：
% \begin{equation}\begin{aligned} ... \\ ... \end{aligned}\end{equation}
% \begin{equation}\begin{split} ... \\ ... \end{split}\end{equation}
% \begin{multline} ... \\ ... \end{multline}

\section{讨论}

正文。

\section{结论}

正文。

\IfFileExists{refs.bib}{\bibliographystyle{unsrt}\bibliography{refs}}{}

\end{document}
